
\documentclass[twocolumn, 10pt]{article}
\usepackage[utf8]{inputenc}
\usepackage[spanish]{babel}
\usepackage{geometry}
\geometry{a4paper, margin=1.5cm}
\usepackage{titlesec}
\usepackage{hyperref}

\title{\textbf{MoneyMatch: Protocolo Descentralizado de Moneymatches para Fighting Games}}
\author{Proyecto Final - Diseño de Soluciones en el Ecosistema Ethereum}
\date{Junio 2026}

\begin{document}

\maketitle

\section{Definición del Problema}
En la comunidad competitiva de juegos de pelea (especialmente en plataformas como Fightcade), los enfrentamientos por dinero o "moneymatches" dependen actualmente de la confianza mutua o de la intervención de intermediarios manuales. Esta estructura presenta ineficiencias críticas: 
\begin{itemize}
    \item \textbf{Riesgo de contraparte:} El perdedor puede negarse a pagar.
    \item \textbf{Dependencia de terceros:} Los árbitros pueden ser lentos, parciales o cobrar comisiones.
    \item \textbf{Falta de registro inmutable:} No existe un historial público y transparente de los retos y sus pruebas (replays) [1].
\end{itemize}

\section{Propuesta de Valor}
FightStake utiliza la infraestructura de \textbf{Ethereum} para actuar como un escrow (garantía) imparcial y automatizado. La solución se basa en tres pilares de Web3 [3, 4]:
\begin{enumerate}
    \item \textbf{Transparencia:} Los fondos están bloqueados en un Smart Contract auditable.
    \item \textbf{Inmutabilidad:} El registro de la partida y su replay quedan grabados de forma permanente.
    \item \textbf{Descentralización:} Se elimina la necesidad de un árbitro central, devolviendo el control a los competidores.
\end{enumerate}

\section{Arquitectura del Sistema}

\subsection{Capa de Smart Contracts}
Se ha diseñado un contrato personalizado en \textbf{Solidity} enfocado en la eficiencia y seguridad [2, 4, 5]:
\begin{itemize}
    \item \textbf{Escrow de Fondos:} El contrato retiene las apuestas de ambos competidores hasta que se cumple una condición de liberación.
    \item \textbf{Registro por Eventos:} Para optimizar el costo de \textbf{gas}, el link del replay no se almacena en el estado del contrato, sino que se emite a través de un \texttt{event}. Esto garantiza que la información sea pública y permanente en los logs de la blockchain de forma económica [4, 5].
    \item \textbf{Mecanismo de Resolución:} El perdedor debe ejecutar una función de "Aceptación de Derrota" para liberar los fondos al ganador.
\end{itemize}

\subsection{Mecanismo de Timelock}
Para evitar el bloqueo indefinido de fondos (deadlock), el contrato implementa un \textbf{Timelock}. Si tras un periodo de 48 horas ninguno de los jugadores confirma el resultado, las funciones de reembolso se habilitan, permitiendo que cada parte recupere su contribución inicial, mitigando fallos económicos [4].

\subsection{Interacción de Usuario}
1. El Retador crea el match depositando fondos y adjuntando el link del replay.
2. El Retado acepta depositando su parte.
3. Tras la partida, el perdedor confirma en la dApp, activando la transferencia automática de la wallet del contrato al ganador [2].

\section{Stack Tecnológico}
Para este MVP se han seleccionado herramientas que priorizan la simplicidad y la seguridad [2, 4]:
\begin{itemize}
    \item \textbf{Lenguaje:} Solidity (Smart Contracts).
    \item \textbf{Entorno:} Remix para compilación y testeo inicial.
    \item \textbf{Red:} Ethereum Testnet (Sepolia) para el despliegue del MVP [6].
    \item \textbf{Librerías de Interfaz:} Ethers.js para la conexión con MetaMask.
\end{itemize}

\section{Análisis de Riesgos}
Siguiendo los criterios de evaluación [4, 6]:
\begin{itemize}
    \item \textbf{Gas:} El uso de eventos reduce drásticamente el costo de registrar replays.
    \item \textbf{Incentivos:} Existe el riesgo de que un perdedor no acepte la derrota por despecho. El Timelock resuelve el bloqueo de fondos, aunque el ganador no cobraría su premio en este escenario (riesgo de adopción).
    \item \textbf{Seguridad:} Al no usar librerías externas como OpenZeppelin, el contrato se mantiene minimalista para reducir la superficie de ataque [4].
\end{itemize}

\end{document}
